%% ma: L12-B08-0157
%% lop: 12
%% chuong: 2
%% bai: 8
%% dang: 12.8.51
%% loai: TN4
%% muc: VD
%% dap_an: "C"
%% nguon: "CĐ2.4.pdf giáo viên gửi 27/10/2026 (M07, Lê Minh Khôi – Nguyễn Anh Khôi, Chương II Chủ đề 4 Các bài toán về tâm tỉ cự), Câu 7"
%% kien_thuc: "Bài 8 (biểu thức tọa độ của phép toán vectơ, độ dài vectơ), Bài 6 (điểm thỏa hệ thức vectơ, quy tắc ba điểm), Bài 7 (hình chiếu lên mặt phẳng tọa độ)"
%% trang_thai: cho-duyet
%% ly_do: "Dạng mới 12.8.51 đề xuất, chờ giáo viên duyệt"
\begin{ex}
Trong không gian $Oxyz$, cho hai điểm $A(3;2;1)$, $B(-2;3;6)$ và điểm $P(x_P;y_P;z_P)$ thay đổi thuộc mặt phẳng $(Oxy)$. Tính giá trị của biểu thức $T=x_P+y_P+z_P$ khi $\left|\overrightarrow{PA}+3\overrightarrow{PB}\right|$ đạt giá trị nhỏ nhất.
\choice{$T=\dfrac{27}4$}{$T=\dfrac{-3}4$}{\True $T=2$}{$T=-2$}
\loigiai{Gọi $I$ thỏa $\overrightarrow{IA}+3\overrightarrow{IB}=\vec 0\Rightarrow I=\dfrac{A+3B}4=\left(-\dfrac34;\dfrac{11}4;\dfrac{19}4\right)$.

$\left|\overrightarrow{PA}+3\overrightarrow{PB}\right|=4PI$ nhỏ nhất khi $P\left(-\dfrac34;\dfrac{11}4;0\right)$. Vậy $T=2$.}
\end{ex}
